\section{Problem Description}
\label{sec:problem-description}

\subsection{Decision setting and modelling scope}

Two complementary models assign a given workload to available capacity.
The aggregate model assigns patients to speciality sessions over several days
and approximates bed demand. The operational model assigns a daily workload
to rooms and evaluates execution in Mesa. They share search mechanisms but
have different decision spaces and objective scales. Static/reactive policy
comparisons concern the operational model.

An operational instance specifies cases, room count, opening hours, turnover
and an optional temporary closure. Rooms and calendars are experimental
assumptions rather than reconstructed hospital resources. All cases are ready
at opening. Days are independent: unstarted cases are counted, without
carryover, new arrivals or clinical prioritisation. The model evaluates a
fixed workload rather than a multi-day waiting list.

\subsection{Cases and available information}

Each case has a local identifier, procedure category and positive predicted
room-occupancy duration. A separate realised duration determines execution
and remains hidden from ordinary scheduling. Historical entry times define
observed occupancy but do not impose new-schedule release times or ordering.
The room baseline uses 2019--2021 procedure medians and a global fallback
for groups with fewer than ten observations; 2022 supplies evaluation cases.
Invalid dates and missing, zero or nonpositive clock intervals are excluded,
without interpreting exit-before-entry as overnight surgery. Durations round
up to integer minutes. Section~\ref{sec:room-design} reports selection counts.

At replanning, visible inputs are statuses, the accepted plan, estimated room
availability and revealed closures. For an ongoing case, availability uses
its predicted finish or at least the next minute if that prediction has
elapsed, followed by turnover. Hidden actual finish is not supplied. The
explicit oracle experiment instead substitutes realised occupancy estimates
while preserving workloads and resources.

\subsection{Aggregate sessions and bed demand}

A patient has speciality, operating minutes and stay days; a session has
speciality, day and minute capacity. The decision assigns each patient to a
session, without a start time or within-session sequence. Search options use
matching speciality when available, otherwise all sessions. This fallback
is computational and does not establish clinical compatibility.

Session load sums assigned durations. Bed occupancy counts the surgery day
through that day plus supplied LOS, inclusively, truncated to the planning
horizon. A zero-day stay therefore occupies one day. The model has a single
aggregate bed capacity, without existing occupants, bed types or explicit
admission/discharge processes. Session excess, bed excess and workload
standard deviation incur soft penalties; returned plans may exceed capacity.
The EDA adapter constructs sessions and derives a speciality proxy from codes,
requiring clinical confirmation. Historical durations are not out-of-sample
predictions unless a separate validated predictor supplies them.

\subsection{Room allocation and operational constraints}

Rooms are interchangeable: the operational model excludes beds, staff,
equipment and speciality eligibility. A solution allocates waiting cases to
rooms. The shared decoder orders each room's cases by increasing predicted
duration, breaking ties by local ID, then computes starts from current time,
opening, estimated availability and known closures. Sequence is fixed by
this rule rather than searched independently.

Every accepted schedule covers each case exactly once as assigned or
explicitly unassigned. Waiting assignments use integer starts and predicted
durations, respect availability and avoid predicted occupancy/turnover
overlap. Cases unable to start strictly before closing become unassigned,
but may be reconsidered at a later replan. Closing prohibits new starts,
not late finishes; required turnover contributes to overtime.

A closure blocks starts from its beginning up to reopening, without
interrupting ongoing non-preemptive work. Reassignment does not change realised
duration. Actual occupancy and turnover guard execution against conflicts
when predictions are optimistic. Planned starts remain lower bounds: cases
do not start early after an unexpectedly short predecessor unless replanned.

\subsection{Disruptions and rescheduling}

Static and reactive policies share one saved initial plan. Static execution
retains room assignments and order, allowing starts to slip when resources
remain occupied or closed. Reactive execution requests a plan at closure and
reopening. The closure is revealed when it begins, including its end time;
duration overruns alone do not trigger replanning.

Only waiting cases may change. Running/completed commitments remain fixed,
and the common validator checks coverage, resources, timing and consistency
with the current state. A failed replan retains the accepted plan; failed
initial planning leaves cases unassigned. Minute steps process completions
and turnover, availability changes, state synchronisation, optional replanning
and eligible starts. Simulated time pauses during solving, so measured
runtime is not itself an imposed patient delay.

\subsection{Objectives and performance criteria}

The room objective prioritises predicted unstarted waiting cases, then
overtime including turnover, then changed waiting decisions. Changes include
room/start revisions and assigned/unassigned transitions; they are inactive
at initial planning. Section~\ref{sec:room-formulation} formalises these
priorities. The aggregate model instead weights session excess, excess
bed-days and imbalance; raw costs cannot be compared across models.

Execution measures completed/unstarted cases, overtime, utilisation, delay
and revisions. Utilisation counts actual occupancy within opening hours,
excluding turnover, divided by configured room-hours; closure does not reduce
that denominator. Overtime sums post-closing occupancy and turnover over
rooms, excluding idle time. Start delay is relative to the initial plan and
measured only for completed cases with an initial assignment, requiring its
own denominator. Revisions count decision changes across replans, potentially
counting a case more than once. Fitness calls, solver time and request counts
measure effort. These definitions support interpretation within the model;
emergency arrivals, clinical urgency, staff constraints and downstream
congestion remain outside its operational scope.
